% c04 —— §5 二次叠 + §6 镜像对称与单值性（开头）（p008–p010 / 原书 p.134–136）
% 承 c03 末「……$\mathcal{D}/\Gamma$ のトロイダル部分コンパクト化が得られる．」句完；
% 本块从原书 p.134 的节标题「5 2 次スタック」起。
% 结尾止于「……$(H^{amb}_{A,\mathbb{Z}},\nabla^{A},\mathscr{F}^{\bullet}_{A},Q_{A})$
% [Iri11, Definition 6.2]。」句完；原书 p.136 末句「局所系 $H^{amb}_{A,\mathbb{Z}}$ の
% ファイバーは……数値」跨页截断，整句归 c05。

% ===== 原书 p.134 =====

\section{二次叠}

设 $\Delta$ 为 $M_{\mathbb{R}}$ 中的自反多面体，$A=\{v_{0}=0,v_{1},\dots,v_{n+r}\}$ 为 $\Delta$ 的格点集合，
考虑\textbf{扇序列}（fan sequence）
\begin{equation}
  0\to\mathbb{L}\to\mathbb{Z}^{n+r}\xrightarrow{\beta}\check{N}\to 0 \tag{5.1}
\end{equation}
这里 $\check{N}=M$，而 $\mathbb{L}$ 是如下同态 $P\colon\mathbb{Z}^{n+r}\to\check{N}$ 的核：它把第 $i$ 个坐标向量
$e_{i}$ 送到 $v_{i}$。
% [Tutor Note: 原书 (5.1) 的箭头标作 $\beta$，正文却把该映射称为 $P$，(5.4) 又用其对偶 $\beta^{\vee}$。此处照原文分别保留 $P$ 与 $\beta$。]
把 $\mathbb{L}$ 的基写作 $\{c^{(p)}\}_{p=1}^{r}$，其中 $c^{(p)}=(c_{1}^{(p)},\dots,c_{n+r}^{(p)})$。
令 $\tilde{v}_{i}=(v_{i},1)\in\check{N}\oplus\mathbb{Z}$，得到正合列
\begin{equation}
  0\to\mathbb{L}\to\mathbb{Z}^{A}\xrightarrow{\tilde{\beta}}\check{N}\oplus\mathbb{Z}\to 0 \tag{5.2}
\end{equation}
这里映射 $\tilde{\beta}$ 把 $e_{i}$ 送到 $\tilde{v}_{i}$，映射 $\mathbb{L}\to\mathbb{Z}^{A}$ 则把 $c^{(p)}\in\mathbb{L}$ 送到
\begin{equation}
  \tilde{c}^{(p)}=(-c_{1}^{(p)}-\cdots-c_{n+r}^{(p)},c_{1}^{(p)},\dots,c_{n+r}^{(p)})\in\mathbb{Z}^{A} \tag{5.3}
\end{equation}
与扇序列 (5.1) 对偶的序列
\begin{equation}
  0\to\check{M}\xrightarrow{\beta^{\vee}}\mathbb{Z}^{n+r}\to\mathbb{L}^{\vee}\to 0 \tag{5.4}
\end{equation}
称为\textbf{除子序列}（divisor sequence）。

对 $\Delta$ 的多面体分割 $\underline{\Delta}=\{\Delta_{1},\dots,\Delta_{k}\}$，令
\begin{equation}
  C(\underline{\Delta})=\{\psi\in\mathbb{R}^{A}\mid\text{对任意 }i=1,\dots,k\text{，}g_{\psi}\text{ 在 }\Delta_{i}\text{ 上都是仿射线性的}\}, \tag{5.5}
\end{equation}
这里 $g_{\psi}\colon\Delta\to\mathbb{R}$ 是与 $\psi\colon A\to\mathbb{R}$ 相伴的分片线性凸函数。
锥 $C(\underline{\Delta})$ 在由 $M_{\mathbb{R}}$ 上的仿射线性函数所成的加法群 $\Aff(M_{\mathbb{R}})$ 的作用下不变，
商锥 $C(\underline{\Delta})/\Aff(M_{\mathbb{R}})$ 确定 $\mathbb{R}^{A}/\Aff(M_{\mathbb{R}})$ 内的完备扇 $\mathcal{F}(A)$。
这个扇 $\mathcal{F}(A)$ 称为\textbf{二次扇}（secondary fan）[GKZ94]。
二次扇 $\mathcal{F}(A)$ 的极大锥对应于 $\Delta$ 的\textbf{正则三角剖分}（coherent triangulation）。
$A$ 的仿射相关子集若其任意真子集都仿射无关，则称为\textbf{回路}（circuit）[GKZ94, 7.1.B]。
相邻的正则三角剖分可通过沿回路的修改互相变换 [GKZ94, Theorem 7.2.10]。

正合列 (5.2) 中映射 $\mathbb{L}\to\mathbb{Z}^{A}$ 的转置 $\mathbb{Z}^{A}\to\mathbb{L}^{\vee}$ 使 $\mathbb{Z}^{A}$ 的基的像构成
$\mathcal{F}(A)$ 的一维锥的生成元，由此在 $\mathcal{F}(A)$ 上定出叠式扇的结构。
这个叠式扇所确定的环面 Deligne–Mumford 叠称为\textbf{二次叠}（secondary stack）[DKK]。

% ===== 原书 p.135 =====

二次叠的稠密环面自然地与 $\mathbb{L}_{\mathbb{C}^{\times}}$ 等同。
二次叠 $X_{\mathcal{F}(A)}$ 的粗模空间是 $\mathbb{P}(\mathbb{C}^{A})$ 关于 $\mathbb{T}$ 的
\textbf{Chow 商}（Chow quotient）[KSZ91, KSZ92]。

对 $\underline{\Delta}$ 的面（即某个 $\Delta_{i}$ 的面）$\Delta'$，令
\[
  C(\underline{\Delta},\Delta')=\{\psi\in C(\underline{\Delta})\mid
  g_{\psi}\text{ 在 }\Delta'\text{ 上取最小值}\},
\]
锥 $C(\underline{\Delta},\Delta')$ 在 $\mathbb{R}\subset\Aff(M_{\mathbb{R}})$ 的加法作用下不变。
各个锥 $C(\underline{\Delta},\Delta')/\mathbb{R}$ 的集合构成 $\mathbb{R}^{A}/\mathbb{R}$ 的完备扇
$\tilde{\mathcal{F}}(A)$。这个扇称为 \textbf{Lafforgue 扇}（Lafforgue fan）[Laf03, Hac]，
与 Lafforgue 扇相伴的环面叠称为 \textbf{Lafforgue 叠}（Lafforgue stack）[DKK]。
自然同态 $\mathbb{Z}^{A}/\mathbb{Z}\to\mathbb{Z}^{A}/\Aff_{\mathbb{Z}}(M)$ 诱导出扇的态射
$\tilde{\mathcal{F}}(A)\to\mathcal{F}(A)$，并定出环面叠的态射
$\varphi_{X}\colon X_{\tilde{\mathcal{F}}(A)}\to X_{\mathcal{F}(A)}$。

二次叠 $X_{\mathcal{F}(A)}$ 的环面轨道对应于 $\Delta$ 的正则多面体分割；
$\varphi_{X}$ 在这个轨道上的纤维，是与该多面体分割对应的环面簇的并。

对 $x=(x_{1},\dots,x_{n})\in(\mathbb{C}^{\times})^{n}$ 与 $a=(a_{0},\dots,a_{n+r})\in\mathbb{C}^{A}$，
Laurent 多项式
\begin{equation}
  W(x;a)=\sum_{i=0}^{n+r}a_{i}x_{1}^{v_{i1}}\cdots x_{n}^{v_{in}} \tag{5.6}
\end{equation}
给出 $(\mathbb{C}^{\times})^{n}\times\mathbb{C}^{A}$ 上的函数。这个函数在 $(\alpha_{i})_{i=1}^{n}\in\mathbb{T}$
对 $(\mathbb{C}^{\times})^{n}\times\mathbb{C}^{A}$ 的作用
\begin{align}
  x_{i}&\mapsto\alpha_{i}x_{i}, & i&=1,\dots,n, \tag{5.7}\\
  a_{k}&\mapsto\alpha_{1}^{-v_{k1}}\cdots\alpha_{n}^{-v_{kn}}a_{k}, & k&=0,\dots,n+r \tag{5.8}
\end{align}
下不变，而在 $\alpha\in\mathbb{C}^{\times}$ 的作用
\begin{align}
  x_{i}&\mapsto x_{i}, & i&=1,\dots,n, \tag{5.9}\\
  a_{k}&\mapsto\alpha a_{k}, & k&=0,\dots,n+r \tag{5.10}
\end{align}
下乘以 $\alpha$。因此 $W(-;a)^{-1}(0)\subset(\mathbb{C}^{\times})^{n}$ 的同构类只依赖于
$[a]\in\mathbb{C}^{A}/(\mathbb{T}\times\mathbb{C}^{\times})$。
对族 $W^{-1}(0)\subset(\mathbb{C}^{\times})^{n}\times\mathbb{C}^{A}\to\mathbb{C}^{A}$
按 $\mathbb{T}\times\mathbb{C}^{\times}$ 的作用取商，并同时紧化底空间与纤维，
便得到超曲面的族 $\varphi_{Y}\colon\mathfrak{Y}_{A}\to X_{\mathcal{F}(A)}$。
换言之，Lafforgue 叠 $X_{\tilde{\mathcal{F}}(A)}$ 上存在自然的线丛
$\mathcal{O}_{X_{\tilde{\mathcal{F}}(A)}}(1)$，它限制到族
$X_{\tilde{\mathcal{F}}(A)}\to X_{\mathcal{F}(A)}$ 的各纤维上给出反典范丛；$W$ 给出它的全局截面，
$\mathfrak{Y}_{A}$ 是它的零点集，而 $\varphi_{Y}$ 是到 $X_{\mathcal{F}(A)}$ 的投影。

Laurent 多项式 $W$ 的临界点集合在 $(\mathbb{C}^{\times})^{n}\times\mathbb{C}^{A}$ 上由
\begin{equation}
  W=\frac{\partial W}{\partial x_{1}}=\cdots=\frac{\partial W}{\partial x_{n}}=0 \tag{5.11}
\end{equation}

% ===== 原书 p.136 =====

定义。它到 $\mathbb{C}^{A}$ 的投影之像的闭包 $\nabla_{A}\subset\mathbb{C}^{A}$ 称为
\textbf{判别式簇}（discriminantal variety），在 $\nabla_{A}$ 上消没的不可约多项式
$\Delta_{A}\in\mathbb{Z}[a_{0},\dots,a_{n+r}]$ 称为 \textbf{$A$ 判别式}（$A$-discriminant），
$\Delta_{A}$ 的 Newton 多面体称为 $A$ 的\textbf{二次多面体}（secondary polytope）。
已知二次多面体的法扇与二次扇一致。
由于 $\nabla_{A}\cap(\mathbb{C}^{\times})^{A}$ 在 $\mathbb{T}\times\mathbb{C}^{\times}$ 的作用下不变，
它下降到 $\mathbb{L}_{\mathbb{C}^{\times}}=(\mathbb{C}^{\times})^{A}/(\mathbb{T}\times\mathbb{C}^{\times})$
中的超曲面 $\overline{\nabla}_{A}$ 上；后者称为
\textbf{既约 $A$ 判别式簇}（reduced $A$-discriminantal variety）。

令
\begin{equation}
  \Phi_{q}(\lambda)=\prod_{j=1}^{n}\Bigl(\sum_{p=1}^{r}c_{j}^{(p)}\lambda_{p}\Bigr)^{c_{j}^{(q)}} \tag{5.12}
\end{equation}
则已知：有理映射 $h\colon\mathbb{P}^{r}\to\mathbb{L}_{\mathbb{C}^{\times}}$，
$\lambda=[\lambda_{1}:\cdots:\lambda_{r+1}]\mapsto(\Phi_{1}(\lambda_{1}),\dots,\Phi_{r}(\lambda))$
的像给出既约 $A$ 判别式簇；这称为既约 $A$ 判别式簇的
\textbf{Horn–Kapranov 单值化}（Horn–Kapranov uniformization）[Hor89, Kap91]。
% [Tutor Note: 原书此处印作 $\Phi_{1}(\lambda_{1})$，其后分量却写作 $\Phi_{r}(\lambda)$。按 (5.12)，$\Phi_{q}$ 依赖于整个 $\lambda$，$\Phi_{1}$ 的自变量 $\lambda_{1}$ 疑为 $\lambda$ 之误。此处照原图录作 $\lambda_{1}$。]

把族 $\varphi_{Y}$ 在判别式簇的补集
$X_{\mathcal{F}(A)}^{reg}=\mathbb{L}_{\mathbb{C}^{\times}}\setminus\nabla_{A}$ 上的限制记作
\begin{equation}
  \varphi_{Y}^{reg}\colon\mathfrak{Y}_{A}^{reg}\to X_{\mathcal{F}(A)}^{reg} \tag{5.13}
\end{equation}

\section{镜像对称与单值性}

由 Hodge 理论，在 $X_{\mathcal{F}(A)}^{reg}$ 上确定出纯极化整 Hodge 结构的变形
$(H_{B,\mathbb{Z}}^{vc},\nabla^{B},\mathscr{F}_{B}^{\bullet},Q_{B})$ [Iri11, Definitions 6.5 and 6.7]。
这里 $H_{B,\mathbb{Z}}^{vc}$ 是 $X_{\mathcal{F}(A)}^{reg}$ 上的局部系统，它在点
$[a]\in X_{\mathcal{F}(A)}^{reg}$ 上的纤维是 $H^{n-1}(Y_{a},\mathbb{Z})$ 中由消没闭链生成的子格；
$\nabla^{B}$ 是 $H_{B,\mathbb{Z}}^{vc}\otimes\mathcal{O}_{X_{\mathcal{F}(A)}^{reg}}$ 上的 Gauss–Manin 联络，
$\mathscr{F}_{B}^{\bullet}$ 是 Hodge 滤过，$Q_{B}$ 是极化。
另一方面，在镜像一侧，把
\begin{equation}
  H^{\bullet}_{amb}(\check{Y};\mathbb{Z})=\Imm\bigl(\iota^{*}\colon
  H^{\bullet}(\check{X};\mathbb{Z})\to H^{\bullet}(\check{Y};\mathbb{Z})\bigr) \tag{6.1}
\end{equation}
取作由环境环面簇的上同调群经限制所得的 $H^{\bullet}(\check{Y};\mathbb{Z})$ 的子空间时，
在以 $H_{amb}^{2}(\check{Y};\mathbb{C})$ 的适当子空间 $U$ 为底空间、以
$H^{\bullet}(\check{Y};\mathbb{Z})$ 为纤维的局部系统 $H_{A,\mathbb{Z}}^{amb}$ 上，
把由 $\check{Y}$ 的 Gromov–Witten 不变量定义的
$\mathscr{H}_{A}=H^{\bullet}_{amb}(\check{Y};\mathbb{C})\otimes\mathcal{O}_{U}$ 上的 Dubrovin 联络 $\nabla^{A}$、
Hodge 滤过 $\mathscr{F}_{A}^{p}=H_{amb}^{4-2p}(\check{Y};\mathbb{C})\otimes\mathcal{O}_{U}$
以及 Poincaré 配对 $Q_{A}\colon\mathscr{H}_{A}\otimes\mathscr{H}_{A}\to\mathcal{O}_{U}$ 组合起来，
便得到纯极化整 Hodge 结构的变形
$(H_{A,\mathbb{Z}}^{amb},\nabla^{A},\mathscr{F}_{A}^{\bullet},Q_{A})$ [Iri11, Definition 6.2]。

% CHUNK c04 DONE pages p008–p010
